\section{Множества}

\begin{definition}{Множество}{}
    \textbf{Множество} --- это совокупность объектов (элементов) произвольной природы,
    обладающих каким-либо общим свойством.
\end{definition}

Запись $x = y$ ($x \neq y$) означает, что $x$ и $y$ --- это один и тот же объект (разные объекты).

Запись $x \in X$ ($x \notin X$) означает, что элемент $x$ принадлежит (не принадлежит) множеству $X$.

\begin{description}
    \item [Пустое множество] $\varnothing$: Множество, которое не содержит ни одного элемента.
    \item [Подмножество] $X \subset Y$: Множество $X$ является подмножеством множества $Y$,
        если каждый элемент из $X$ также лежит в $Y$.
    \item [Равенство] $X = Y$: Множества равны, если $X \subset Y$ и $Y \subset X$.
    \item [Объединение] $X \cup Y$: Объединение множеств состоит из элементов,
        которые принадлежат хотя бы одному из этих множеств.
    \item [Пересечение] $X \cap Y$: Пересечение множеств состоит из элементов,
        которые принадлежат обоим множествам одновременно.
    \item [Разность] $X \setminus Y$: Разность множеств состоит из элементов,
        которые принадлежат $X$ и при этом не принадлежат $Y$.
\end{description}

\begin{definition}{Декартово произведение}{}
    Множество всех возможных упорядоченных пар элементов $(x, y)$, где $x \in X$, а $y \in Y$,
    называется \textbf{декартовым произведением} множеств $X \times Y$.
\end{definition}

\begin{definition}{Бинарное отношение}{}
    \textbf{Бинарное отношение} на множестве $X$ задается подмножеством $R \subset X \times X$.
\end{definition}

Запись $x_1 R x_2$ означает, что элементы $x_1, x_2 \in X$ связаны бинарным отношением $R$.

\begin{definition}{Отношение эквивалентности}{}
    Бинарное отношение называется \textbf{отношением эквивалентности},
    если оно обладает следующими свойствами:
    \begin{enumerate}
        \item \textit{Рефлексивность:} Для любого $x \in X$ верно $x R x$.
        \item \textit{Симметрично:} Если $x_1 R x_2$, то $x_2 R x_1$.
        \item \textit{Транзитивно:} Если $x_1 R x_2$ и $x_2 R x_3$, то $x_1 R x_3$.
    \end{enumerate}
\end{definition}

\begin{theorem}{Классы эквивалентности}{}
    Отношение эквивалентности $R$ разбивает все множество $X$
    на попарно непересекающиеся подмножества (\textbf{классы эквивалентности}) по правилу:
    два элемента $x_1$ и $x_2$ принадлежат одному классу в том и только в том случае, когда $x_1 R x_2$.
\end{theorem}
